\section{Ordinary steps on the given space}\label{sec:absolute}

We construct the contraction and flip of every negative analytic ray in
arbitrary dimension.  The input is an ordinary rational klt pair in the
\emph{global Weil-divisor} category.  Generalized pairs enter only through
the cone and base-point-free theorems; the steps themselves are ordinary.
The main point is to preserve actual line bundles, including their fixed
Cartier indices, rather than only numerical classes.

\subsection{Integral descent}

Bundles are written additively in numerical expressions.  For example,
$L-(K_T+C)$ denotes the rational line bundle obtained from $L$ and the
rational adjoint.  All degrees in the following lemma are degrees on
compact curves contracted by the indicated morphism.

\begin{lemma}[Integral descent]\label{lem:descent}
Let $f:T\to Z$ be a projective bimeromorphic contraction of normal complex
spaces, and let $L$ be a line bundle numerically trivial over $Z$.
Suppose that every point of $Z$ has a neighbourhood $U$ on which there is
an ordinary rational klt pair $(f^{-1}U,C_U)$ with rational adjoint and
\[
 L-(K_{f^{-1}U}+C_U)\quad\text{$f$-nef over $U$}.
\]
Then $f_*L$ is a line bundle and the evaluation map is an isomorphism
\begin{equation}\label{eq:integral-descent}
 f^*f_*L\simeq L.
\end{equation}
In particular the conclusion applies if the local ordinary adjoints are
relatively antiample.  It also applies to a global ordinary rational klt
pair $(T,C)$ whenever $L-(K_T+C)$ is $f$-nef.  Clearing one denominator
therefore descends a rational line bundle as an actual rational line
bundle.
\end{lemma}

\begin{proof}
Fix $z\in Z$ and shrink to a connected Stein neighbourhood $U$ on which
the stated pair is defined.  A line bundle $M$ on $f^{-1}U$ has a
Cartier-divisor representative there: the coherent direct image $f_*M$
has rank one over the nonempty open where $f$ is an isomorphism, so
Cartan's Theorem~A supplies a nonzero section.  Its pullback is a section
of $M$ that is not identically zero, and its zero divisor is Cartier.
This argument takes place over $U$, not on the whole compact space.

Every rational line bundle $N$ on $f^{-1}U$ is also relatively big.
Indeed, choose $q>0$ with $qN$ a line bundle and an $f$-ample line bundle
$A$.  The same direct-image argument gives a nonzero section of
$qN-A$, hence
\[
 N\simq q^{-1}A+q^{-1}E\qquad(E\geq0).
\]
This is relative bigness in the analytic sense
\cite[Definition~2.46]{Fujino2022}.  Apply it to
$N=L-(K_T+C_U)$.  The klt base-point-free theorem gives, after shrinking
around $z$, relative generation of $L^k$ for \emph{every} sufficiently
large integer $k$ \cite[Theorem~6.2 and Remark~6.3]{Fujino2022}.

Choose finitely many such generating sections near the compact fibre
$F=f^{-1}(z)$.  The morphism they define on the projective reduced fibre
$F_{\mathrm{red}}$ is constant on each irreducible component: a
positive-dimensional image would yield a curve with positive degree
against the pullback of the hyperplane bundle, contrary to the
numerical triviality of $L^k$.  Such a curve is obtained by successive
general hyperplane sections of a component.  Connectedness of $F$ makes
the images of all its components the same point.  A linear combination
of the generating sections is consequently nonzero at every point of
$F_{\mathrm{red}}$.

This section is a generator of $L^k$ in the local ring at every point of
$F$: its residue is nonzero, and hence its representing function is a
unit.  This also treats a nonreduced fibre.  Properness allows us to
shrink $U$ away from the image of the section's zero set.  Thus $L^k$ is
trivial on $f^{-1}U$.  Apply the argument to consecutive integers $k$ and
$k+1$, and use one common neighbourhood.  Their quotient trivializes
$L$ itself.  Since $f_*\cO_T=\cO_Z$, on this neighbourhood $f_*L\simeq
\cO_U$, and evaluation is an isomorphism.  These local conclusions give
\eqref{eq:integral-descent} globally.  No additional tensor power enters
the integral conclusion.
\end{proof}

The nef variant is important even when the adjoint has degree zero over
$f$.  For example, if $P=K_T+C$ is a rational klt adjoint, $\ell P$ is
Cartier, and $P$ is numerically trivial over $f$, apply
Lemma~\ref{lem:descent} to $L=\ell P$.  Its difference from the adjoint
is $(\ell-1)P$, so the \emph{same} index $\ell$ descends.

\subsection{Supporting classes and projectivity}

A class is \emph{modified big} if it is the birational pushforward of a big class
\cite[Definition~2.8]{HX2026}. In particular, a big class is modified big.

We use the following consequences of the cone and transcendental
base-point-free theorems \cite[Theorems~1.3 and~1.4]{HX2026}.  For a
compact K\"ahler gklt pair with a descending nef part $\beta$, the cone
$\NA(T)$ is the sum of its adjoint-nonnegative part and rays of rational
curves.  Only finitely many rays are needed when the boundary plus
$\beta$ is big.  If its adjoint class $\alpha$ is nef and the boundary
plus $\beta$ is modified big, there is a Moishezon contraction
$f:T\to Z$ to a normal compact K\"ahler space with $\alpha=f^*\gamma$
for a K\"ahler class $\gamma$ on $Z$.  Neither assertion requires strong
factoriality.  Projectivity in our weaker category will be proved from
the ordinary rational adjoint.

\begin{lemma}[Supports with a K\"ahler margin]\label{lem:absolute-support}
Let $(T,B)$ be an ordinary rational klt pair on a normal compact K\"ahler
space.  Assume that $B$ is $\Q$-Cartier and that $D=K_T+B$ is a rational
line bundle.  If $D$ is not nef, $\NA(T)$ has a $D$-negative extremal ray.
Every such ray $R$ is generated by a rational curve and has a nonempty
relatively open set $\mathcal U_R\subset R^\perp$ of classes satisfying
\begin{equation}\label{eq:absolute-support}
 \alpha\text{ is nef},\qquad
 \NA(T)\cap\alpha^\perp=R,\qquad
 \alpha-c_1(D)\text{ is K\"ahler}.
\end{equation}
\end{lemma}

\begin{proof}
Fix a K\"ahler class $\omega$ and write $d=c_1(D)$.  Cone duality applies
because klt singularities are rational.  The slice
\[
 S=\{v\in\NA(T):\omega\cdot v=1\}
\]
is compact.  Indeed, fix smooth representatives of a basis of $\BC(T)$.
Each is bounded above and below by a multiple of a K\"ahler form, so all
its pairings on $S$ are bounded; the slice is closed.  If $D$ is not nef,
the minimum of $d$ on $S$ is negative.  An extreme point of its minimum
face gives a negative extremal ray.

Fix any such ray and let $r\in S$ be its normalized point.  Choose
$\varepsilon>0$ with $(d+\varepsilon\omega)\cdot r<0$.  The descending
nef part $\varepsilon\omega$ changes no discrepancy, and
$B+\varepsilon\omega$ is big since $B$ is effective and $\Q$-Cartier.
The finite form of the cone theorem gives
\begin{equation}\label{eq:absolute-finite-cone}
 \NA(T)=\NA(T)_{d+\varepsilon\omega\geq0}
       +\sum_{j=1}^{N}\R_{\geq0}[C_j],
\end{equation}
where the $C_j$ are rational curves.  Normalize their classes in $S$
and discard repetitions.  By extremality, $r$ is one of these points,
say $r_1$.  Let $K$ be the convex hull of the others and of
$S\cap\{d+\varepsilon\omega\geq0\}$.  This is compact and excludes
$r$; otherwise extremality of $r$ would fail.

If $K$ is nonempty, strict separation gives a linear functional $h$
with $h(r)=0$ and $h>0$ on $K$, after subtracting a multiple of $\omega$.
Regard it as a Bott--Chern class by duality.  For all sufficiently small
$\delta>0$, $h-\delta d$ is strictly positive on $K$ and at $r$.
Since $S=\operatorname{conv}(\{r\}\cup K)$, it is strictly positive on
$S$, hence K\"ahler.  Thus $\alpha=h/\delta$ has the required properties.
Its positive minimum on $K$ persists under small perturbations in
$R^\perp$; openness of the K\"ahler cone also preserves
$\alpha-d$.  These perturbations give $\mathcal U_R$.  If $K$ is empty,
then $S=\{r\}$ and $-d$ is K\"ahler; a small neighbourhood of $0$ in
$R^\perp$ has the required properties.  This includes the case
$R^\perp=\{0\}$.
\end{proof}

\begin{lemma}[Relative positivity]\label{lem:relative-positivity}
Let $f:T\to Z$ be a proper morphism of compact complex spaces, with $T$
and $Z$ K\"ahler.  Suppose that $L$ is a rational line bundle and
\[
 c_1(L)=\kappa-f^*\gamma,
\]
where $\kappa$ is K\"ahler on $T$ and $\gamma$ is a smooth real
Bott--Chern class on $Z$.  Then $L$ is $f$-ample and $f$ is projective.
Conversely, a normal space projective over a compact K\"ahler space is
compact K\"ahler if it has a global relatively ample line bundle.
\end{lemma}

\begin{proof}
Clear the denominator of $L$ and choose any smooth Hermitian metric
on that line bundle. The Bott--Chern identity says that the difference
between its curvature and the indicated difference of forms is
$\ddc$ of a global smooth function. Adjusting the metric by the
exponential of that function gives the required curvature, with the
same denominator restored at the end. Smooth functions and metrics
here have their normal-space meaning in local embeddings.
On every fibre its curvature is positive.  Positivity of a line bundle
on a compact complex space implies ampleness, and for a proper analytic
map ampleness on every fibre is equivalent to relative ampleness
\cite[Corollary~1.12 and Definition~3.1--Remark~3.2]{FujinoVanishing2026}.
This criterion does not presuppose projectivity and applies to
nonreduced fibres as well.  Existence of the relatively ample line bundle
makes $f$ projective.

For the converse, local relative embeddings give metrics on a relatively
ample bundle with positive curvature in fibre directions.  Glue their
weights by a partition of unity on the base.  The resulting metric
retains that positivity.  Add a sufficiently large multiple of a base
K\"ahler form to its curvature.  The sum is strictly positive; compactness
allows one multiple to work everywhere.
\end{proof}

\subsection{The ordinary contraction and flip}

We can now perform each step and preserve the category.  The next
proposition records both the numerical ranks of the negative contraction
and the actual canonical algebra of a flip.  This description identifies the transformed adjoint as an actual
rational line bundle.

\begin{proposition}[Ordinary negative steps]\label{prop:ordinary-step}
Let $(T,B)$ be an ordinary rational klt pair on a normal compact K\"ahler
space, globally $\Q$-factorial in the Weil-divisor sense, with canonical
sheaf a rational line bundle.  Put $D=K_T+B$.  Every $D$-negative
extremal ray $R\subset\NA(T)$ has a projective contraction $f:T\to Z$
with connected fibres and normal compact K\"ahler target, such that
\begin{equation}\label{eq:ordinary-negative-bc}
 f^*\BC(Z)=R^\perp,\qquad
 \rho(T/Z)=1,\qquad -D\text{ is }f\text{-ample}.
\end{equation}
Here $\rho$ is the rank of global line-bundle degrees on contracted
curves.  The target has rational singularities and
$R^j f_*\cO_T=0$ for $j>0$.  If $\dim Z<\dim T$, this is a Mori fibre
space.  Otherwise $f$ is bimeromorphic, and precisely one of the
following occurs.
\begin{enumerate}[label=\textup{(\roman*)}]
\item The exceptional locus is one prime divisor $E$.  Put $T'=Z$,
$B'=f_*B$ and $D'=K_{T'}+B'$.  Then
\begin{equation}\label{eq:ordinary-divisorial-difference}
 D=f^*D'+eE\qquad(e>0).
\end{equation}
\item The contraction is small.  For an adjoint Cartier index $r>0$,
the intrinsic algebra
\begin{equation}\label{eq:ordinary-flip-algebra}
 \cR_r=\bigoplus_{m\geq0}f_*\cO_T(mrD)
\end{equation}
is locally finitely generated, and $T'=\Proj_Z\cR_r$ is the ordinary
flip.  It has a small projective morphism $f^+:T'\to Z$ with connected
fibres.  Put $B'=\phi_*B$ and $D'=K_{T'}+B'$, where
$\phi:T\dashrightarrow T'$ is the induced small map.  For a sufficiently
divisible choice of $r$ there is an actual isomorphism
\begin{equation}\label{eq:ordinary-tautological}
 \cO_{\Proj_Z\cR_r}(1)\simeq\cO_{T'}(rD').
\end{equation}
In particular $D'$ is a rational line bundle and is $f^+$-ample.
Passing to a sufficiently divisible Veronese leaves $T'$ unchanged.
\end{enumerate}
In both birational cases, $T'$ is compact K\"ahler and globally
Weil-divisor $\Q$-factorial, its canonical sheaf is a rational line
bundle, and $(T',B')$ is klt.  Strong global $\Q$-factoriality is
preserved when imposed on $T$.  A supplied global canonical Weil divisor
is transported to a canonical Weil divisor on $T'$.  No prime divisor is
extracted.  All log discrepancies weakly increase; the increase is strict
for every original prime contracted by the step, and, for a flip, for
every place centred in either exceptional locus.
\end{proposition}

\begin{proof}
Choose $\alpha\in\mathcal U_R$ from Lemma~\ref{lem:absolute-support}
and put $\kappa=\alpha-c_1(D)$.  Insert the descending nef part
$\kappa$ into the ordinary pair.  It changes no discrepancy, while
$B+\kappa$ is big.  Transcendental base-point-freeness gives
$f:T\to Z$ and $\gamma$ K\"ahler with $\alpha=f^*\gamma$.
Thus $c_1(-D)=\kappa-f^*\gamma$, and
Lemma~\ref{lem:relative-positivity} proves projectivity and relative
ampleness.  This establishes projectivity without strong factoriality.
For every curve $C\subset T$,
\begin{equation}\label{eq:ordinary-contracted-curves}
 f(C)\text{ is a point}\quad\Longleftrightarrow\quad
 \alpha\cdot C=0\quad\Longleftrightarrow\quad[C]\in R.
\end{equation}
The rational curve generating $R$ is contracted, so $f$ is nontrivial.

Relative vanishing gives $R^j f_*\cO_T=0$ for $j>0$
\cite[Theorem~5.2]{Fujino2022}; the log Fano base has rational
singularities \cite[Lemma~8.8]{DHP2024}.  Lemma~\ref{lem:pullback-image}
now identifies $f^*\BC(Z)$ with the classes of degree zero on all
vertical curves, which is exactly $R^\perp$.  The global line-bundle
rank is also one, since all those curves are proportional and the
rational line bundle $D$ detects their ray.  This proves
\eqref{eq:ordinary-negative-bc}.  The lower-dimensional case has all
the stated Mori fibre-space properties.  In equal dimension,
connectedness gives generic degree one, so $f$ is bimeromorphic.

\smallskip\noindent
\emph{Divisorial contractions.}
Suppose $E$ is an exceptional prime.  It is $\Q$-Cartier.  If
$E\cdot R\geq0$, it would be $f$-nef, contradicting negativity for a
nonzero effective exceptional divisor.  Thus $E\cdot R<0$.  Every
vertical curve lies in $E$, since an effective $\Q$-Cartier divisor
has nonnegative degree on a curve outside its support.  Every
positive-dimensional projective fibre is covered by curves.  A
zero-dimensional local component of a connected fibre would be isolated,
and the birational form of Zariski's main theorem would make $f$ an
isomorphism there.  Consequently $\Exc(f)=E$.

For a prime Weil divisor $A_Z$ on $Z$, let $A_T$ be its strict transform.
Choose $t\in\Q$ so $(A_T+tE)\cdot R=0$.  The coefficient is rational
by taking the ratio of the two rational degrees on one integral curve.
A Cartier multiple of $A_T+tE$ descends by Lemma~\ref{lem:descent}:
its difference from $D$ is relatively ample.  The descended line bundle
agrees with the corresponding reflexive power of $\cO_Z(A_Z)$ off
$f(E)$, a subset of codimension at least two.  Reflexivity identifies
them everywhere.  Hence every global prime on $Z$ is $\Q$-Cartier.

If a global canonical Weil divisor is given, push it forward.  It agrees
with the canonical sheaf on the big open where $f$ is an isomorphism, so
reflexivity shows that it remains a canonical Weil representative.
Alternatively, start with an invertible reflexive power of $\omega_T$,
adjust it by a rational multiple of $E$ to degree zero, and descend after
clearing denominators.  On the same big open the result is a reflexive
power of $\omega_Z$; hence that power is invertible everywhere.  This
proves the canonical assertion without assuming a global meromorphic
canonical section.

Now $D'$ is a rational line bundle and its canonical identification away
from $E$ gives \eqref{eq:ordinary-divisorial-difference}.  Intersection
with $R$ gives $e>0$.  If $b_E$ is the coefficient of $E$ in $B$, then
\[
 a(E;Z,B')=1-b_E+e>1-b_E=a(E;T,B)>0.
\]
Pullback of the effective divisor $eE$ gives weak increase for all other
log discrepancies.  Thus the target pair is klt.

\smallskip\noindent
\emph{The small canonical model.}
Suppose $f$ is small.  Finite generation of the ordinary relative log
canonical ring \cite[Theorem~1.18]{Fujino2022} gives local finite
generation of \eqref{eq:ordinary-flip-algebra}.  When canonical data are
written as sheaves, this application can be made over each Stein open
in $Z$: nonzero sections of rank-one proper direct images provide the
meromorphic representatives there.  Changing representatives induces
compatible graded isomorphisms of the intrinsic algebra.

The ordinary analytic flip theorem \cite[Theorem~1.14]{Fujino2022}
constructs precisely its relative canonical model.  It has no local
factoriality hypothesis.  Its proof glues the unique local log canonical
models, giving a normal small projective $f^+:T'\to Z$ globally.  The
fibres are connected because $Z$ is normal.  A finite cover of the
compact base permits one common sufficiently divisible Veronese
generated in degree one.  Replace $r$ by that multiple.  The resulting
tautological bundle is relatively ample and agrees with
$\cO_T(rD)$ on the common big open.  Its reflexive extension is
$\cO_{T'}(rD')$, proving \eqref{eq:ordinary-tautological} as an actual
bundle isomorphism.  This also proves that $D'$ is a rational line
bundle.  The relative Proj is invariant under this Veronese operation.
By Lemma~\ref{lem:relative-positivity}, $T'$ is compact K\"ahler.

For a global prime $A'$ on $T'$, transform it to $A$ on $T$ and choose
$t\in\Q$ so $(A+tD)\cdot R=0$.  Descend a Cartier multiple by
Lemma~\ref{lem:descent}.  On $T'$ its pullback agrees on the common big
open with the same multiple of $A'+tD'$.  Reflexivity extends the
identification; since $D'$ is already a rational line bundle, $A'$ is
$\Q$-Cartier.  This proves global Weil-divisor factoriality.  A given
canonical Weil divisor is transported on that common open and then
reflexively.  In the sheaf formulation, repeat the descent argument with
an invertible power of $\omega_T$ adjusted by $tD$; undoing the twist
by $tD'$ proves the rational-line-bundle property of $\omega_{T'}$.

On a common smooth resolution $p:V\to T$, $q:V\to T'$,
Lemma~\ref{lem:negativity} gives the actual discrepancy divisor
\begin{equation}\label{eq:ordinary-small-difference}
 p^*D=q^*D'+F,\qquad F\geq0,
\end{equation}
exceptional over both sides, with the asserted strictness over the
flipping and flipped loci.  This also proves klt preservation.

\smallskip\noindent
\emph{Optional strong factoriality.}
In the divisorial case, take any coherent rank-one reflexive sheaf
$\cF$ on $Z$.  Its reflexive pullback to $T$ is coherent of rank one.
Strong factoriality supplies an invertible reflexive power.  Adjust it
by a rational multiple of $E$, descend as above, and compare off $f(E)$.
A reflexive power of $\cF$ is therefore a line bundle.

In the small case start with such a sheaf $\cF'$ on $T'$.  On a common
smooth resolution put
\[
 \cM=(q^*\cF'/\text{torsion})^{**},\qquad
 \cF=(p_*\cM)^{**}.
\]
Here $\cM$ is a line bundle and $\cF$ is coherent reflexive of rank one;
both give the required transform on the common big open.  An invertible
power of $\cF$, adjusted by a rational multiple of $D$, descends.
Pull back the descended bundle and undo the twist by $D'$.  After
clearing denominators, reflexivity identifies the resulting line bundle
with a reflexive power of $\cF'$.  Thus the strong condition is
preserved.  Only sheaves defined on the whole compact spaces have been
used.

\end{proof}

\subsection{Transport of Bott--Chern classes}

For later scaling we need to transport one fixed nef datum.  The following
construction uses the rank of the negative contraction only.  It makes
no assertion that all Bott--Chern classes on a flipped space come from
the preceding space.

\begin{lemma}[Class transport]\label{lem:class-transport}
For a birational step of Proposition~\ref{prop:ordinary-step}, let
$f':T'\to Z$ mean $f^+$ in the small case and $\id_Z$ in the divisorial
case.  Every $\theta\in\BC(T)$ has a unique expression
$\theta=f^*\eta+s\,c_1(D)$, with $\eta\in\BC(Z)$ and $s\in\R$.
Define
\begin{equation}\label{eq:class-transport}
 \theta'= (f')^*\eta+s\,c_1(D').
\end{equation}
On a common resolution, $p^*\theta-q^*\theta'$ is the class of an
actual divisor exceptional over $T'$.  This rule sends the class of a
global Weil divisor to that of its transform or pushforward.  Moreover,
it propagates the fixed-nef-part trace relation of
Lemma~\ref{lem:class-trace}, with an actual effective exceptional
correction on the new model.  The same conclusions hold for ordinary
real-boundary flips on strongly globally $\Q$-factorial compact K\"ahler
spaces constructed by \cite[Proposition~4.2]{HX2026}.
\end{lemma}

\begin{proof}
Since $D\cdot R\ne0$, subtracting a unique multiple of $c_1(D)$ makes
$\theta$ vanish on $R$.  Equation~\eqref{eq:ordinary-negative-bc} and
pullback injectivity give the unique $\eta$.  On a common resolution
over $Z$ the two pullbacks of $\eta$ agree.  Therefore
\[
 p^*\theta-q^*\theta'
       =s\,[p^*D-q^*D'].
\]
The expression in brackets is the actual exceptional discrepancy
divisor already constructed in the proof of the proposition.

If $\theta=c_1(A)$ for a global Weil divisor $A$, the difference
$p^*A-q^*A'$ is likewise exceptional over $T'$, where $A'$ is its
transform or pushforward.  Subtracting the two identities represents
$q^*(\theta'-c_1(A'))$ by a $q$-exceptional divisor.  That divisor is
numerically trivial over $q$, so negativity in both signs makes it
zero.  Pullback injectivity gives $\theta'=c_1(A')$.

Finally suppose the old trace $H_T$ of a fixed nef class on an earlier
model has, on a common high resolution, the relation
\[
 p^*H_T=p_0^*H+[J_T].
\]
If $H_T=f^*\eta+s\,c_1(D)$, the new correction is the actual divisor
\[
 J_{T'}=J_T-s(p^*D-q^*D'),\qquad
 q^*H_{T'}=p_0^*H+[J_{T'}].
\]
Every old exceptional divisor remains exceptional
over $T'$: neither kind of step extracts divisors.  The resulting
correction $J_{T'}$ is therefore an actual $q$-exceptional divisor.
Its negative is $q$-nef because $p_0^*H$ is nef.  Negativity gives
$J_{T'}\geq0$, and uniqueness and compatibility are those of
Lemma~\ref{lem:class-trace}.

For the stated real-boundary extension, the cited contraction theorem
supplies the negative-side Bott--Chern rank and the ordinary real flip.
The difference $p^*D-q^*D'$ is still an actual exceptional real
discrepancy divisor, by Lemma~\ref{lem:negativity}.  The decomposition,
transport formula, and correction argument above therefore apply
unchanged.  Rationality of $D$ was not used in these arguments.
\end{proof}

\begin{remark}\label{rem:ordinary-continuation}
Proposition~\ref{prop:ordinary-step} starts on the given pair and can be
reapplied after every finite prefix of ordinary steps.  It allows every
negative ray.  Its construction is independent of any termination
claim, and makes no assertion about the positive-side quotient
$\BC(T')/(f^+)^*\BC(Z)$ or about equality of the full Bott--Chern ranks
across a flip.
\end{remark}
